
\chapter{Midterm Review and Sample Questions}
\label{ch:midterm-review-and-sample-questions}
\section*{Midterm Overview}
\addcontentsline{toc}{section}{Midterm Overview}

The midterm is intended to evaluate the student’s understanding of the
foundational concepts developed in the first half of the course. Coverage
typically includes Chapters 1--5 together with the accompanying classroom
discussion and examples.

\subsection*{Recommended format}
\addcontentsline{toc}{subsection}{Recommended format}
\begin{itemize}
  \item \textbf{Part A:} short-answer and definition questions on DSS concepts,
  architecture, and decision-making terminology.
  \item \textbf{Part B:} analytical questions that compare models, evaluation
  approaches, and deployment choices.
  \item \textbf{Part C:} one short essay or case-based design question requiring
  a justified DSS design decision.
\end{itemize}

\subsection*{Preparation guidance}
\addcontentsline{toc}{subsection}{Preparation guidance}
\begin{itemize}
  \item Revise key distinctions such as MIS vs.\ BI vs.\ DSS and structured vs.\
  semi-structured vs.\ unstructured decisions.
  \item Be able to explain bounded rationality, common cognitive biases, and why
  governance matters in DSS.
  \item Review model management vocabulary including verification, validation,
  calibration, and data leakage.
  \item Practice describing how a prediction becomes an operational decision.
\end{itemize}

\section*{Sample Questions}
\addcontentsline{toc}{section}{Sample Questions}

\begin{enumerate}[label=\textbf{Q\arabic*.}]
  \item \textbf{Define a DSS and explain how it differs from a traditional MIS.}

  \textit{Brief answer guideline:} A strong answer should mention the focus on
  semi-structured and unstructured decisions, the role of models and interactive
  analysis, and the difference between reporting and decision support.

  \item \textbf{Why is bounded rationality important when designing a DSS
  interface?}

  \textit{Brief answer guideline:} A strong answer should connect bounded
  rationality to limited time, limited cognition, progressive disclosure,
  summaries, and bias mitigation.

  \item \textbf{Explain the difference between verification and validation in
  model management.}

  \textit{Brief answer guideline:} Verification asks whether the model was built
  correctly; validation asks whether the model is appropriate for the actual
  decision context.

  \item \textbf{A hospital wants to use an ML model to prioritize emergency
  cases. Which evaluation issues matter more than accuracy alone?}

  \textit{Brief answer guideline:} Expected points include recall for critical
  cases, calibration, threshold selection, operational capacity, fairness, and
  human review.
\end{enumerate}
